\subsection{Concentration of three-cycle parameters}
\label{sec:three-cycle-concentration}

The two-cycle estimate on a compact positive source interval also
locates the coefficient sequences that can carry three cycles. We use
the canonical slice throughout: for a field \(F\) near \(F_*\), put
\(\theta_F=\theta(\mathscr C_0(F))\), with
\(\mathscr C_0\) from \eqref{eq:fixed-coefficient-transport} and
\(\theta\) from Lemma~\ref{lem:slice}. Define
\(\gamma(F),\alpha_0(F),\alpha_1(F)\) by
\eqref{eq:center-parameters} at \(\theta_F\), and, when
\(\alpha_0(F)<0\), set
\[
 s(F)=\sqrt{-\alpha_0(F)},\qquad
 p(F)=\frac{\alpha_1(F)}{s(F)}.
\]
Thus the ratio below refers to this fixed analytic slice.

\begin{proposition}[Uniform two-cycle exclusion on compact positive ratios]
\label{prop:compact-positive-two-cycle-exclusion}
For every fixed \(0<\eta\le P<\infty\), there are an open
coefficient neighborhood \(\mathcal V_{\eta,P}\) of \(F_*\) in
the full twelve-coefficient quadratic germ and
\(\eps_{\eta,P}>0\) such that every
\(F\in\mathcal V_{\eta,P}\) with
\(\alpha_0(F)<0\) and \(\eta\le p(F)\le P\) has at most
two distinct isolated periodic orbits \(C\) satisfying
\[
 d_{\rm H}(\iota_*(C),\Gamma_{I4})<\eps_{\eta,P}.
\]
The return estimate holds on one complete physical receiver, retaining
the original windows and clips and all inputs between its roots.
\end{proposition}

\begin{proof}
Apply Lemma~\ref{ann:lem-compact-positive-domains} and
Proposition~\ref{sharp:upper-annulus} on the enlarged fixed ratio
interval \([\eta/2,2P]\). They give the two-zero multiplicity bound
on one complete physical receiver, including its endpoint, middle and
high ranges. Choose its cap so that
\[
 0<b_{\eta,P}\le\min\{h_c,z_*\},\qquad
 L_{\rm raw}b_{\eta,P}<d_{\rm out}/4,
 \qquad |\alpha_0|<d_{\rm out}/4.
\]
All constants may depend on \(\eta,P\). The ordinary coefficient
restriction in \eqref{eq:finite-multiplier} makes the actual finite
multiplier satisfy the final bound used by that proposition.

Now choose
\(0<h_{\rm in}<\min\{b_{\eta,P}/2,d_{\rm in}/(4L_{\rm raw})\}\)
and \(|\alpha_0|<d_{\rm in}/4\). The height estimate
\eqref{eq:height-capture} and Lemma~\ref{lem:loc-capture} admit every
sufficiently close whole cycle to this counted receiver, at its unique
prime section hit. Lemma~\ref{lem:slice-bound-transfer}, with
\(\mathcal B=\{\alpha_0<0,\eta\le p\le P\}\) and
\(N_{\rm cyc}=2\), gives the full-germ neighborhood and collar.
\end{proof}

\begin{corollary}[Three-cycle concentration]
\label{cor:three-cycle-concentration}
Let \(F_n\to F_*\) in the full quadratic coefficient germ. Suppose
each \(F_n\) has three distinct isolated periodic orbits
\(C_n^1,C_n^2,C_n^3\), all of this same field, and
\[
 \max_{1\le j\le3}
 d_{\rm H}(\iota_*(C_n^j),\Gamma_{I4})\longrightarrow0.
\]
Then \(\alpha_0(F_n)<0\) and \(\alpha_1(F_n)>0\) for all
sufficiently large \(n\), and
\[
 p(F_n)=\frac{\alpha_1(F_n)}{\sqrt{-\alpha_0(F_n)}}
 \longrightarrow0.
\]
\end{corollary}

\begin{proof}
The whole-coefficient bounds of
Proposition~\ref{sharp:upper-coefficient-bounds}, on the common
receiver followed by geometric capture, give at most two cycles when
\(\alpha_0\ge0\), at most one when
\(\alpha_0<0,\alpha_1\le0\), and at most two when
\(\alpha_0<0,p\ge P_*\), for one fixed \(P_*<\infty\).
In the last range the exact trace scaling is
\(r_L=\alpha_1\), \(p_L=1\), \(a_L=-p^{-2}\);
its independent small \(r_L,|\gamma|\) bounds and positive unscaled
cap cover the entire large-ratio sector.
The same-field transport used above puts these exclusions on one full
coefficient neighborhood and collar. The sequence and all three cycle
sequences eventually enter them. Hence eventually
\(\alpha_0(F_n)<0\), \(\alpha_1(F_n)>0\), and
\(0<p(F_n)<P_*\).

Fix \(0<\eta<P_*\). Apply
Proposition~\ref{prop:compact-positive-two-cycle-exclusion} to
\([\eta,2P_*]\). Its neighborhood and collar are fixed for this
\(\eta\), so the sequence and its three cycles eventually enter
them. If \(p(F_n)\ge\eta\) there, all three cycles belong to
one two-cycle sector, a contradiction. Thus \(p(F_n)<\eta\)
eventually. This proves the limit. The argument uses one fixed
\(\eta\) at a time and requires no common collar as
\(\eta\downarrow0\).
\end{proof}

\begin{corollary}[A consequence in a positive-gamma subregion]
\label{cor:positive-gamma-concentration}
Along any subsequence in Corollary~\ref{cor:three-cycle-concentration}
for which \(\gamma(F_n)>0\), if
\[
 \frac{s(F_n)}{\gamma(F_n)^{3/2}}\longrightarrow0,
\]
then
\[
 \liminf\frac{p(F_n)}{\sqrt{\gamma(F_n)}}\ge\sqrt{12}.
\]
\end{corollary}

\begin{proof}
Fix \(0<\pi_0<\sqrt{12}\). The second alternative of
Lemma~\ref{deg:lem-raw-jacobian} supplies a fixed \(k_0>0\)
and a uniformly contracting complete return when
\[
 \frac{s}{\gamma^{3/2}}\le k_0,\qquad
 \frac{p}{\sqrt\gamma}\le\pi_0.
\]
Its fixed-window endpoint quotient bounds the derivative on the
whole physical interval, including intervening inputs. After its
coefficient and multiplier restriction, geometric capture gives one
neighborhood and collar. The subsequence eventually satisfies the
first inequality and enters both; three isolated cycles exclude the
second. Thus its lower limit is at least every
\(\pi_0<\sqrt{12}\), which proves the assertion.
\end{proof}
